% Generated draft: requires mathematical and source review.
\documentclass[11pt]{article}
\usepackage[margin=25mm]{geometry}
\usepackage{fontspec}
\setmainfont{Latin Modern Roman}
\usepackage{amsmath,amssymb,mathtools,unicode-math}
\setmathfont{Latin Modern Math}
\usepackage{enumitem,needspace,xurl,hyperref}
\setlength{\parindent}{0pt}
\setlength{\parskip}{5pt}
\setlength{\emergencystretch}{3em}
\newcommand{\N}{\mathbb{N}}
\newcommand{\Z}{\mathbb{Z}}
\newcommand{\Q}{\mathbb{Q}}
\newcommand{\R}{\mathbb{R}}
\newcommand{\C}{\mathbb{C}}
\newcommand{\F}{\mathbb{F}}
\newcommand{\A}{\mathbb{A}}
\newcommand{\PP}{\mathbb{P}}
\newcommand{\E}{\mathbb{E}}
\newcommand{\eps}{\varepsilon}
\newcommand{\dd}{\,\mathrm{d}}
\begin{document}
\section*{UnsolvedMath ID: 30005123}
\textbf{Source problem code:} catalogue entry
\subsection{Definitions and mathematical statement}
For orientation, the displayed constructions can be described in the standard finite-quiver setting over \(\C\). This setting makes the formulas below meaningful; the supplied record does not explicitly specify the permitted class of quivers, so it must not be read as a recovered hypothesis of the original question.

Let \(Q=(Q_0,Q_1)\) be a finite quiver. Its doubled quiver \(\overline Q\) is obtained by adjoining an opposite arrow \(a^*\) for each \(a\in Q_1\). Writing \([x,y]=xy-yx\), its preprojective algebra is
\[
\Pi_Q=\C\overline Q\Big/\left\langle\sum_{a\in Q_1}[a,a^*]\right\rangle.
\]
The tripled quiver \(\widetilde Q\) is obtained from \(\overline Q\) by adjoining a loop \(\omega_i\) at every vertex \(i\in Q_0\). Its cubic potential is the cyclic class
\[
\widetilde W=
\left(\sum_{a\in Q_1}[a,a^*]\right)
\left(\sum_{i\in Q_0}\omega_i\right)
\quad\text{in }\C\widetilde Q/[\C\widetilde Q,\C\widetilde Q].
\]
Here the commutator space in the quotient is a vector subspace. More generally, a potential on a quiver \(R\) is a linear combination of cyclic paths, considered modulo commutators. For a cyclic path \(b_1\cdots b_\ell\), define
\[
\partial_b(b_1\cdots b_\ell)
 =\sum_{\substack{1\leq j\leq\ell\\b_j=b}}
 b_{j+1}\cdots b_\ell b_1\cdots b_{j-1},
\]
and extend linearly. The associated Jacobi algebra is
\[
\operatorname{Jac}(R,W)=
\C R/\langle\partial_bW\mid b\in R_1\rangle.
\]
The supplied excerpt states that
\[
\operatorname{Jac}(\widetilde Q,\widetilde W)
\cong\Pi_Q[\omega],
\qquad
\omega\longleftrightarrow\sum_{i\in Q_0}\omega_i,
\]
where \(\omega\) is central.

Write \(A_{\Pi_Q}\) for the cohomological Hall algebra associated with the preprojective algebra, and \(A_{R,W}\) for the cohomological Hall algebra associated with a quiver with potential in the source's construction. These symbols retain the source notation: the supplied excerpt does not give the complete cohomology, vanishing-cycle, grading, or multiplication conventions needed to define these algebras precisely.

\textbf{Research question.} Can the preprojective cohomological Hall-algebra construction be modified to produce a fermionized, or partially odd, version recovering Yangians of classical Lie superalgebras, with \(\mathfrak{gl}(m\mid n)\) as the stated example? In the quiver-with-potential formulation, can one construct a quiver \(Q'\) and a potential \(W'\) such that \(A_{Q',W'}\) is a partial fermionization of \(A_{\widetilde Q,\widetilde W}\), together with a precise identification of the resulting algebra with the intended Yangian or intended part of it?

For nonnegative integers \(m,n\), the notation \(\mathfrak{gl}(m\mid n)\) denotes the Lie superalgebra of endomorphisms of \(\C^m\oplus\C^n\), with the first summand even, the second odd, and bracket
\[
[x,y]=xy-(-1)^{|x||y|}yx
\]
for homogeneous endomorphisms. This explains the example without imposing a quantifier absent from the source: the record does not specify which pairs \((m,n)\) are required, which classical Lie superalgebras are included, or how their parameters determine \(Q\) and \((Q',W')\).

\textbf{Unresolved qualifications.} The expression ``partial fermionization'' is not defined by a mathematical condition in the supplied record. The excerpt motivates it through parity changes and the addition of quadratic terms to potentials, but does not specify which terms, which graded pieces change parity, or which structures must be preserved. It also does not specify the presentation, deformation parameters, completion, or precise portion of the target Yangian. The clean statement calls the two questions equivalent; the supplied excerpt gives related questions without enough definitions to establish that equivalence. Consequently, the question above records a construction target, not a fully specified existence predicate.

The supplied contextual discussion describes the original algebra as concentrated in even cohomological degrees and suggests an interpretation as one half of a generalized Yangian, saying that this can be made precise in finite and affine type. It also records, as work in progress with L. Hennecart and S. Schlegel-Mejia, a generalized Kac--Moody interpretation of a Lie algebra denoted \(\mathfrak g_{\Pi_Q}\). These are source-reported context, not additional conclusions established here. In particular, the distinction between recovering an entire Yangian and recovering one half remains unresolved.
\subsection{Short English statement}
The question asks whether changing a quiver and its potential can introduce odd parity into the cohomological Hall algebra of a preprojective algebra and thereby recover Yangians associated with classical Lie superalgebras, such as \(\mathfrak{gl}(m\mid n)\). The source does not define partial fermionization or specify the exact Yangian identification and quantifiers, so these remain explicit formulation gaps. The supplied dated curation labels the problem partially solved; this draft does not independently verify that assessment.
\subsection{Sources}
\begin{itemize}
\item Supplied local UnsolvedMath registry record, permanent ID \(30005123\), title ``Fermionized Yangians from Quivers with Potential,'' collection ``Oberwolfach Reports Open Problems,'' exact source problem code \texttt{catalogue entry}. No proposer is named in the record; the proposed year is recorded as 2022.
\item Supplied report citation: \emph{Interactions between Algebraic Geometry and Noncommutative Algebra (2023)}, Oberwolfach Reports, pp.\ 1239--1312, DOI \texttt{10.4171/owr/2022/23}: \url{https://doi.org/10.4171/owr/2022/23}. This is the authoritative report reference supplied by the registry; no exact problem page or speaker attribution is supplied.
\item Supplied literature reference, labelled ``Davison, fermionization paper'': \url{https://doi.org/10.1017/S001708952200009X}. The record does not supply its full bibliographic citation.
\item Supplied literature reference, labelled ``Oberwolfach report'': \url{https://doi.org/10.4171/owr/2022/23}.
\item Literature assessment dated exactly \texttt{2026-08-22}: the supplied curation labels the problem ``partially solved'' and reports that Davison constructs a parity-switching fermionization for preprojective cohomological Hall algebras and relates it to Yangian phenomena, while the full requested realization of all classical Lie-superalgebra Yangians is not established. This is supplied dated curation, not independently verified literature assessment; its time-sensitive status and cited claims require independent checking.
\item The registry labels the statement \texttt{corrected\_verified} and reports comparison with the official Oberwolfach report. These are supplied verification metadata, not verification performed for this draft. The excerpt attributes related work in progress to collaboration with L. Hennecart and S. Schlegel-Mejia, without identifying the speaker in this record.
\end{itemize}
\end{document}
